\documentclass{../../PublicResources/DocClassLight}

    \DocumentTitle{文档标题测试}
    \DocumentCreatedDate{2020/7/21}

    \AuthorName{Mr. Kin}

\begin{document}
    \pagenumbering{Roman} % 大写罗马字母样式页码。
    \maketitle
    \addcontentsline{toc}{chapter}{封面}
    \frontmatter
    {\centering \tableofcontents} % 生成目录页。
    \mainmatter

    % 正文
    \part{中文字体测试}
    \chapter{章标题测试}
    \section{节标题测试}
    \subsection{子节标题测试}
    \subsubsection{子子标题测试}
    \paragraph{段标题测试}
    测试
    \subparagraph{子段标题测试}
    测试

    \subsection{字体样式测试}
    {\rmfamily 思宋：测试 \textbf{粗体：测试} \textup{直立体：测试} \textit{意大利斜体：测试} \textsl{倾斜体：测试}}

    {\sffamily 思黑：测试 \textbf{粗体：测试} \textup{直立体：测试} \textit{意大利斜体：测试} \textsl{倾斜体：测试}}

    \subsection{字体测试}
    {\rmfamily 思宋：滚滚长江东逝水，浪花淘尽英雄。是非成败转头空，青山依旧在，几度夕阳红。白发渔樵江渚上，惯看秋月春风。一壶浊酒喜相逢，古今多少事，都付笑谈中。}

    {\sffamily 思黑：滚滚长江东逝水，浪花淘尽英雄。是非成败转头空，青山依旧在，几度夕阳红。白发渔樵江渚上，惯看秋月春风。一壶浊酒喜相逢，古今多少事，都付笑谈中。}

    \part{English Font Test}
    \chapter{Chapter Title Test}
    \section{Section Title Test}
    \subsection{Subsection Title Test}
    \subsubsection{Subsubsection Title Test}
    \paragraph{Paragraph Title Test}
    Test
    \subparagraph{Subparagraph Title Test}
    Test

    \subsection{Font Format Test}
    {\rmfamily textrm: test \textbf{textbf：test} \textup{textup：test} \textit{textit：test} \textsl{textsl：test}}

    {\sffamily textsf: test \textbf{textbf：test} \textup{textup：test} \textit{textit：test} \textsl{textsl：test}}

    {\ttfamily texttt: test \textbf{textbf：test} \textup{textup：test} \textit{textit：test} \textsl{textsl：test}}

    \subsection{Font Test}
    {\rmfamily textrm: The quick brown fox jumps over the lazy dog}

    {\sffamily textsf: The quick brown fox jumps over the lazy dog}

    {\ttfamily texttt: The quick brown fox jumps over the lazy dog}

    \subsection{IPA Test}
    anode \ipa{'ænoʊd}

    cathode \ipa{'kæθoʊd}

    \chapter{各类测试}

    \begin{intro}
        intro环境测试。intro 环境，用以章节开头的文字介绍。排版上比普通正文多缩进两个文字。
    \end{intro}

    \paragraph{普通文字段} 滚滚长江东逝水，浪花淘尽英雄。是非成败转头空，青山依旧在，几度夕阳红。白发渔樵江渚上，惯看秋月春风。一壶浊酒喜相逢，古今多少事，都付笑谈中。

    \paragraph{有序列表}
    \begin{enumerate}[label={Step \arabic*.}]
        \item test
        \item test
        \item test
        \item 滚滚长江东逝水，浪花淘尽英雄。是非成败转头空，青山依旧在，几度夕阳红。白发渔樵江渚上，惯看秋月春风。一壶浊酒喜相逢，古今多少事，都付笑谈中。
    \end{enumerate}

    \paragraph{无序列表}
    \begin{itemize}
        \item test
        \item test
        \item test
        \item 滚滚长江东逝水，浪花淘尽英雄。是非成败转头空，青山依旧在，几度夕阳红。白发渔樵江渚上，惯看秋月春风。一壶浊酒喜相逢，古今多少事，都付笑谈中。
    \end{itemize}

    \begin{lstlisting}[language={C},title={\textsf{C语言代码段测试}}]
        for(,1,);
        printf("Hello World!"); // 注释测试 Test
    \end{lstlisting}

    \paragraph{\color{black}颜色测试}
    {\color{red} 红色}

    \paragraph{下划线测试}
    \uline{下划线}
    \uuline{双下划线}
    \uwave{波浪线}
    \sout{删除线}
    \xout{斜线}
    \dashuline{下划线-虚线}
    \dotuline{下划线-点}

    \rmnum{3}

    \Rmnum{3} $K_{test}^{\text{\Rmnum{3}}}$ $K_{test}^{\mathrm{\Rmnum{3}}}$

    \newpage
    \noindent {\bfseries \sffamily 图片测试}
    \begin{figure}[h!]
        \centering
        \includegraphics[scale=0.9]{SampleImage}
        \caption{示例图片}
    \end{figure}

    \nocite{*} % 不使用 cite 也能生成参考文献。
    \printbibliography % 生成参考文献排版。
    \addcontentsline{toc}{chapter}{参考文献} % 添加参考文献进目录。

    \appendix
    % 附录
    \chapter{附录测试1}
    \section{节测试}
    \subsection{子节测试}
    \chapter{附录测试2}

\end{document}
